\section{Introduction}
\label{sec:intro}

Video datasets underpin advances in computer vision, but their publication raises privacy concerns when subjects have not consented to automated tracking. SAM2~\citep{ravi2024sam2} has emerged as a high-accuracy, promptable video segmentation and tracking system that third parties can readily use to extract subject trajectories from published data. A natural defense is a \emph{publisher-side video preprocessor}: a system that adds visually imperceptible perturbations to video frames before release, causing SAM2 to fail on any queried object while preserving the utility of the video for other research tasks.

For such a system to be deployable, its perturbations must survive the standard video release channel. Video datasets are almost universally compressed with H.264 (or H.265) codec before distribution---YouTube, Vimeo, and most institutional repositories apply CRF 23 by default, achieving 38--42 dB PSNR while discarding high-frequency spatial content. This codec step is not optional: it is the mechanism by which large video corpora are made practically distributable.

The central question of this paper is therefore: \textbf{can pixel-constrained adversarial perturbations survive H.264 CRF23 and disrupt SAM2 video tracking?}

We conducted seven experiments across two attack paradigms and four candidate insertion points to answer this question. Within the $\varepsilon = 8/255$ L$_\infty$ budget, the answer is \textbf{no}. Every tested strategy---from end-to-end learned pixel-level CNN preprocessors to feature-space direct optimization targeting SAM2's memory encoder and object pointer tokens---produced near-zero tracking disruption after CRF23 compression.

\paragraph{Contributions.}
\begin{enumerate}
  \item We present the first systematic study of \emph{codec robustness} for adversarial attacks against SAM2 video tracking, covering pixel-level and feature-space strategies across four attack insertion points.
  \item We document a striking separation: attacks that devastate SAM2 tracking without codec ($\djfadv$ up to $+0.87$) fail completely after CRF23 ($\djfauc \approx 0$). We term this the \emph{codec-kill} phenomenon.
  \item We hypothesize and discuss the root cause as the high-frequency concentration of $L_\infty$-constrained perturbations and their interaction with DCT quantization in H.264, and identify the key open question: whether BPDA/EOT-style codec-aware optimization can overcome this barrier.
  \item We establish that, for the tested direct-optimization attacks, the \emph{optimizer is not the bottleneck}: the per-video Adam optimization finds effective perturbations (e.g., $\djfadv = +0.867$ on the bus video); the codec then destroys them. This provides evidence against the hypothesis that stronger gradient-based optimization alone would recover codec robustness within $\varepsilon = 8/255$.
\end{enumerate}

\paragraph{Why this matters.}
Prior adversarial attack papers against video segmentation~\citep{uap_sam2_2025,darksam2024} evaluate under standard (codec-free) conditions. Our finding suggests these results may not transfer to real deployment pipelines. Adversarial robustness benchmarks and privacy tool evaluations in video settings should include codec compression as a standard evaluation condition.

\paragraph{Paper organization.}
\cref{sec:background} covers SAM2 architecture, H.264 compression, and the threat model. \cref{sec:related} reviews related work. \cref{sec:attacks} describes both attack paradigms. \cref{sec:pixel_exps} presents pixel-level results. \cref{sec:feature_exps} presents feature-space results. \cref{sec:analysis} analyzes the codec-kill mechanism. \cref{sec:discussion} discusses implications and limitations.
